% Part: computability
% Chapter: computability-theory
% Section: fixed-point-thm

\documentclass[../../../include/open-logic-section]{subfiles}

\begin{document}

\olfileid{cmp}{thy}{fix}
\olsection{د ثابت ټکي قضيه}

راځئ يو ځل بيا د درېدنې مسئله وګورو۔ د لنډمهالې نښې په توګه
~$\gn{\cfind{x}(y)}$ د~$\tuple{x, y}$ دپاره وليکو؛ دا د
$\cfind{x}(y)$ د قيمت د يو ``نوم'' استازيتوب وګڼئ۔ په دې نښه د
درېدنې د مسئلې د ناپرېکړتيا يو ثبوت بيا بيانولے شو۔

پوښتنه: ايا داسې محاسبه کېدونکې تابع~$h$ شته چې لاندې خاصيت ولري؟ د
هر~$x$ او~$y$ دپاره
\[
h(\gn{\cfind{x}(y)}) =
\begin{cases}
1 & \text{که $\cfind{x}(y) \fdefined$ وي} \\
0 & \text{که داسې نۀ وي۔}
\end{cases}
\]

ځواب: نه؛ که نه، دا جزوي تابع
\[
g(x) \simeq
\begin{cases}
0 & \text{که $h(\gn{\cfind{x}(x)}) = 0$ وي} \\
\text{تعريف شوې نۀ ده} & \text{که داسې نۀ وي}
\end{cases}
\]
به محاسبه کېدونکې وه، او ځکه به ئې کوم شاخص~$e$ درلود۔ خو بيا لرو:
\[
\cfind{e}(e) \simeq
\begin{cases}
0 & \text{که $h(\gn{\cfind{e}(e)}) = 0$ وي} \\
\text{تعريف شوې نۀ ده} & \text{که داسې نۀ وي،}
\end{cases}
\]
چې په دې صورت کښې~$\cfind{e}(e)$ هله او يوازې هله تعريف شوې ده چې
تعريف شوې نۀ وي، او دا تناقض دے۔

اوس هغه معادله وګورئ چې~$\cfind{e}$ لري۔ دلته په يوه معنا د
ځان-مراجعې يوه بېلګه شته: موږ د~$\cfind{e}(e)$ قيمت داسې برابر کړے
چې په يوه ځانګړې طريقه پر~$\gn{\cfind{e}(e)}$ تکيه وکړي۔ د ثابت ټکي
قضيه وايي چې دا کار په عمومي ډول \emph{کولے شو}---يوازې د تناقضونو د
ثابتولو دپاره نۀ۔

\olref{lem:fixed-equiv} د ثابت ټکي د قضيې د بيانولو دوه معادل شکلونه
راکوي۔ د منطقي نظره، د دې بيانونو معادلتوب له دې راځي چې دواړه رښتيا
دي؛ خو زموږ اصلي موخه دا ده چې هر يو په نېغه توګه له بل څخه راځي، نو
د همدې قضيې د بديلو بيانونو په توګه اخيستل کېدے شي۔

\begin{lem}
\ollabel{lem:fixed-equiv}
لاندې بيانونه معادل دي:
\begin{enumerate}
\item د هرې جزوي محاسبه کېدونکې تابع~$g(x,y)$ دپاره داسې شاخص~$e$
  شته چې د هر~$y$ دپاره
\[
\cfind{e}(y) \simeq g(e,y).
\]
\item د هرې محاسبه کېدونکې تابع~$f(x)$ دپاره داسې شاخص~$e$ شته چې
  د هر~$y$ دپاره
\[
\cfind{e}(y) \simeq \cfind{f(e)}(y).
\]
\end{enumerate}
\end{lem}

\begin{proof}
$(1) \Rightarrow (2)$: ورکړې~$f$ ته~$g$ داسې تعريف کړئ چې
$g(x,y) \simeq \fn{Un}(f(x),y)$ وي۔ (1) وکاروئ څو داسې شاخص~$e$
تر لاسه کړئ چې د هر~$y$ دپاره
\begin{align*}
\cfind{e}(y) & \simeq \fn{Un}(f(e),y) \\
& \simeq \cfind{f(e)}(y).
\end{align*}

$(2) \Rightarrow (1)$: ورکړې~$g$ ته د~$s$-$m$-$n$ قضيه وکاروئ څو
داسې~$f$ تر لاسه کړئ چې د هر~$x$ او~$y$ دپاره
$\cfind{f(x)}(y) \simeq g(x,y)$ وي۔ (2) وکاروئ څو داسې شاخص~$e$ تر
لاسه کړئ چې
\begin{align*}
\cfind{e}(y) & \simeq \cfind{f(e)}(y) \\
& \simeq g(e,y).
\end{align*}
په دې سره ثبوت بشپړېږي۔
\textbf{د ژباړې د نسخې سمون (OLCMP-038):} سرچينې په دې دوو
استنتاجي لړيو کښې د جزوي تابعو د قيمتونو تر منځ د عادي برابرۍ نښه
کارولې وه، سره له دې چې اړخونه تعريف شوي نۀ هم کېدے شي۔ دلته د همدې
څپرکي ټاکلې جزوي برابرۍ نښه کارول شوې ده، څو دواړه تعريف شوي او
ناتعريف حالتونه په کره ډول پوښل شي۔
\end{proof}

\begin{explain}
مخکښې له دې چې وښيو بيان~(1) رښتيا دے (او ځکه~(2) هم)، فکر وکړئ چې
دا څومره عجيبه خبره ده۔ $e$ يو کمپيوټري پروګرام وګڼئ؛ بيان~(1) وايي
چې د هرې ورکړې جزوي محاسبه کېدونکې~$g(x,y)$ دپاره داسې کمپيوټري
پروګرام~$e$ موندلے شئ چې $g_e(y) \simeq g(e,y)$ محاسبه کوي۔ په بله
وينا، داسې کمپيوټري پروګرام موندلے شئ چې يوه پر خپل پروګرام مراجعه
کوونکې تابع محاسبه کوي۔
\end{explain}

\begin{thm}
په~\olref{lem:fixed-equiv} کښې دواړه بيانونه رښتيا دي۔ په ځانګړي
ډول، د هرې جزوي محاسبه کېدونکې تابع~$g(x,y)$ دپاره داسې شاخص~$e$
شته چې د هر~$y$ دپاره
\[
\cfind{e}(y) \simeq g(e,y).
\]
\end{thm}

\begin{proof}
اړين توکي لا د پورته درېدنې مسئلې په بحث کښې نغښتي دي۔
$\fn{diag}(x)$ داسې محاسبه کېدونکې تابع وي چې د هر~$x$ دپاره د تابع
$f_x(y) \simeq \cfind{x}(x,y)$ يو شاخص راګرځوي، يعنې
\[
\cfind{\fn{diag}(x)}(y) \simeq \cfind{x}(x,y).
\]
$\fn{diag}$ داسې تابع وګڼئ چې د يوې دوه‌ځاييزې تابع پروګرام د يوې
يوځاييزې تابع په پروګرام بدلوي، او دا د اصلي پروګرام په لومړي دليل
کښې د هغه خپل شاخص په ثابتولو تر لاسه کېږي۔ تابع~$\fn{diag}$ په صوري
ډول داسې تعريفېدے شي: لومړے~$s$ داسې تعريف کړئ:
\[
s(x,y) \simeq \fn{Un}^2(x,x,y),
\]
چېرته چې~$\fn{Un}^2$ يوه درې‌ځاييزه تابع ده چې د جزوي محاسبه
کېدونکو دوه‌ځاييزو تابعو دپاره نړيواله ده۔ بيا د~$s$-$m$-$n$ قضيې
له مخې داسې بنسټيزه بازګشتي تابع~$\fn{diag}$ موندلے شو چې دا پوره
کړي:
\[
\cfind{\fn{diag}(x)}(y) \simeq s(x,y).
\]

اوس تابع~$l$ داسې تعريف کړئ:
\[
l(x,y) \simeq g(\fn{diag}(x),y).
\]
او~$\gn{l}$ د~$l$ يو شاخص وي۔ په پای کښې
$e = \fn{diag}(\gn{l})$ وټاکئ۔ بيا د هر~$y$ دپاره لرو:
\begin{align*}
\cfind{e}(y) & \simeq \cfind{\fn{diag}(\gn{l})}(y) \\
& \simeq \cfind{\gn{l}}(\gn{l}, y) \\
& \simeq l(\gn{l}, y) \\
& \simeq g(\fn{diag}(\gn{l}),y) \\
& \simeq g(e, y),
\end{align*}
لکه چې غوښتل شوي وو۔
\end{proof}

\begin{explain}
دلته څه روان دي؟ فرض کړئ تاسو ته د داسې کمپيوټري پروګرام د ليکلو
دنده درکړل شوې چې خپل متن چاپ کړي۔ بيا فرض کړئ چې په داسې پروګرامي
ژبه کار کوئ چې د رشتو يو بډای او عجيب کتابتون لري۔ په ځانګړي ډول،
فرض کړئ ستاسو پروګرامي ژبه داسې تابع~$\fn{diag}$ لري: تابع
$\fn{diag}$ يوه ننوت رشته~$s$ اخلي، په~$s$ کښې د `x' نښې هره بېلګه
مومي، او هغه د اصلي
رشتې په اقتباس شوې بڼه بدلوي۔ د بېلګې په توګه، که دا رشته
\begin{quote}
\begin{verbatim}
hello x world
\end{verbatim}
\end{quote}
د ننوت په توګه ورکړل شي، تابع دا
\begin{quote}
\begin{verbatim}
hello 'hello x world' world
\end{verbatim}
\end{quote}
د وت په توګه راګرځوي۔ په دې صورت کښې موخه شوے پروګرام ليکل اسان
دي؛ ازمويلے شئ چې
\begin{quote}
\begin{verbatim}
print(diag('print(diag(x))'))
\end{verbatim}
\end{quote}
کار سر ته رسوي۔ د~C++ او جاوا په شان ډېرو عامو پروګرامي ژبو کښې هم
همدغه نظر (په يو څه پېچلې پليينې سره) کار کوي۔

موږ د ثابت ټکي د قضيې له ثبوت څخه يوازې څو ګامه لرې يو۔ فرض کړئ د
چاپ تابعې يوه بڼه~$\fn{print}(x,y)$ يوه رشته~$x$ او بل عددي دليل~$y$
اخلي، او رشته~$x$، $y$ ځله چاپ کوي۔ بيا دا ``پروګرام''
\begin{quote}
\begin{verbatim}
getinput(y);
print(diag('getinput(y); print(diag(x), y)'), y)
\end{verbatim}
\end{quote}
پر ننوت~$y$ خپل ځان~$y$ ځله چاپ کوي۔ که د
$\fn{getinput}$---$\fn{print}$---$\fn{diag}$ چوکاټ په هرې تابع
$g(x,y)$ بدل کړو، دا تر لاسه کېږي:
\begin{quote}
\begin{verbatim}
g(diag('g(diag(x), y)'), y)
\end{verbatim}
\end{quote}
دا داسې پروګرام دے چې پر ننوت~$y$، تابع~$g$ پر خپل پروګرام او~$y$
چلوي۔ که ``اقتباس کول'' د ``يو شاخص کارولو'' په توګه وګڼو، پورته
ثبوت تر لاسه کوو۔

اوس دپاره که غواړئ ثبوت صوري چل، يا توره جادوګري، وګڼئ پروا نۀ کوي۔
خو بايد د پورته ورکړل شوي استدلال جزئيات بيا جوړ کړے شئ۔ کله چې د
نابشپړتيا قضيې (او اړونده ``د ثابت ټکي قضيه'') ثابتوو، د دې د کار
کولو د علت د پوهېدو په نورو لارو به بحث وکړو۔
\end{explain}

\begin{tagblock}{lambda}
\begin{digress}
همدغه نظر د ``ثابت ټکي'' يو ترکيبګر هم راکولے شي۔ فرض کړئ يو لامبډا
ترم~$g$ لرئ، او داسې بل ترم~$k$ غواړئ چې~$k$ له~$gk$ سره
$\beta$-معادل وي۔ زموږ د نښيزو دودونو په کارولو ترمونه داسې تعريف
کړئ:
\[
\fn{diag}(x) = xx
\]
او
\[
l(x) = g(\fn{diag}(x))
\]
يعنې~$l$ دا ترم دے: $\lambd[x][g(xx)]$۔ $k$ ترم~$ll$ وي۔ بيا لرو:
\begin{align*}
k & = (\lambd[x][g(xx)])(\lambd[x][g(xx)]) \\
& \red  g((\lambd[x][g(xx)])(\lambd[x][g(xx)])) \\
& = gk.
\end{align*}
که
\[
Y = \lambd[g][((\lambd[x][g(xx)])(\lambd[x][g(xx)]))]
\]
واخلو، نو~$Yg$ او~$g(Yg)$ يو ګډ ترم ته راکمېږي؛ ځکه
$Yg \equiv_\beta g(Yg)$۔ دې ته ``د کري ترکيبګر'' ويل کېږي۔ که پر
ځاے ئې
\[
Y = (\lambd[xg][g(xxg)])(\lambd[xg][g(xxg)])
\]
واخلو، نو په حقيقت کښې~$Yg$، $g(Yg)$ ته راکمېږي، چې لا پياوړے بيان
دے۔ د~$Y$ دا وروستۍ بڼه ``د ټيورينګ ترکيبګر'' بلل کېږي۔
\end{digress}
\end{tagblock}

\end{document}
